% Chapter 4, Section 1 _Linear Algebra_ Jim Hefferon
%  http://joshua.smcvt.edu/linearalgebra
%  2001-Jun-12
\chapter{Keserupaan}
Kita telah menunjukkan bahwa bagi sebarang homomorfisme terdapat basis $B$
dan~$D$ sehingga matriks yang merepresentasikan pemetaan tersebut memiliki
bentuk blok identitas parsial.
\begin{equation*}
  \rep{h}{B,D}
  =
    \begin{pmat}{c|c}
       \text{\textit{Identitas}}  &\text{\textit{Nol}}   \\
       \hline
       \text{\textit{Nol}}      &\text{\textit{Nol}}
    \end{pmat}
\end{equation*}
Representasi ini mendeskripsikan pemetaan yang membawa
\( c_1\vec{\beta}_1+\dots+c_n\vec{\beta}_n \) ke
\(c_1\vec{\delta}_1+\dots+c_k\vec{\delta}_k+\zero+\dots+\zero \), dengan $n$
sebagai dimensi domain dan \( k \) sebagai dimensi daerah hasil.
Dalam representasi ini, tindakan pemetaan mudah dipahami karena sebagian besar
entri matriks bernilai nol.

Bab ini meninjau kasus khusus ketika domain dan kodomain sama.
Dalam hal ini, wajar jika kita menghendaki basis domain dan basis kodomain juga
sama.
Artinya, kita menginginkan basis \( B \) sehingga \( \rep{t}{B,B} \) sesederhana
mungkin, dengan `sederhana' berarti memiliki banyak entri nol.
Kita akan mendapati bahwa matriks berbentuk blok identitas parsial seperti di
atas tidak selalu dapat diperoleh, tetapi kita akan mengembangkan bentuk yang
mendekatinya, yaitu representasi yang hampir diagonal.













\section{Ruang Vektor Kompleks}
\index{ruang vektor!di atas bilangan kompleks}
%<*ReasonForShiftToComplexNumbers>
Dalam bab ini kita perlu memfaktorkan polinomial.
Namun, banyak polinomial tidak dapat difaktorkan di atas bilangan real.
Sebagai contoh, \( x^2+1 \) tidak dapat difaktorkan menjadi hasil kali dua
polinomial linear berkoefisien real; untuk itu diperlukan bilangan kompleks,
$x^2+1=(x-i)(x+i)$.

Karena itu, dalam bab ini kita akan menggunakan bilangan kompleks sebagai
skalar, termasuk sebagai entri vektor dan matriks.
Artinya, kita beralih dari mempelajari ruang vektor di atas bilangan real ke
ruang vektor di atas bilangan kompleks.\index{bilangan kompleks!ruang vektor di atas}
Setiap bilangan real merupakan bilangan kompleks dan sebagian besar contoh
dalam bab ini hanya menggunakan bilangan real. Meskipun demikian, teorema
pentingnya mengharuskan skalar berupa bilangan kompleks.
%</ReasonForShiftToComplexNumbers>
Jadi, bagian pertama ini merupakan tinjauan bilangan kompleks.

Dalam buku ini, kita beralih ke konteks yang lebih umum dengan menggunakan
skalar kompleks karena alasan pragmatis:~peralihan tersebut diperlukan untuk
melanjutkan pembahasan.
Namun, gagasan membentuk ruang vektor dengan mengambil skalar dari struktur
selain bilangan real juga menarik dan berguna.
Penyajian menarik yang menggunakan pendekatan ini sejak awal terdapat dalam
\cite{Halmos} dan \cite{HoffmanKunze}.






\subsectionoptional{Faktorisasi Polinomial dan Bilangan Kompleks}
\textit{Subbagian ini hanya merupakan tinjauan.
  Untuk pembahasan lengkap, termasuk pembuktian, lihat \cite{Ebbinghaus}.}

%<*PolynomialReview0>
Tinjau suatu polinomial\index{polinomial}
$p(x)=c_nx^n+\dots+c_1x+c_0$ dengan
koefisien utama\index{polinomial!koefisien utama}
$c_n\neq 0$.
Kita menyebutnya polinomial \definend{berderajat}~$n$
\index{polinomial!derajat}\index{derajat polinomial}.
Jika $n=0$, maka $p$ merupakan polinomial konstan
\index{polinomial!konstan}\index{polinomial konstan}, $p(x)=c_0$.
Polinomial konstan yang bukan polinomial nol, yaitu $c_0\neq 0$, berderajat
nol.
Kita definisikan polinomial nol berderajat $-\infty$.
%</PolynomialReview0>

\begin{remark}
Mendefinisikan derajat polinomial nol sebagai $-\infty$ %, which most authors do, 
membuat persamaan $\text{derajat}(fg)=\text{derajat}(f)+\text{derajat}(g)$
berlaku bagi semua polinomial.
\end{remark}

%<*PolynomialReview1>
Seperti bilangan bulat, polinomial juga memiliki operasi pembagian; misalnya,
`\( 4 \) termuat \( 5 \) kali dalam \( 21 \) dengan sisa \( 1 \)'.
%</PolynomialReview1>

\begin{theorem}[Teorema Pembagian bagi Polinomial] \label{th:EuclidForPolys}
\index{teorema pembagian}\index{polinomial!teorema pembagian}
%<*th:EuclidForPolys>
Misalkan \( p(x) \) suatu polinomial.
Jika \( d(x) \) merupakan polinomial tak nol, maka terdapat polinomial
\definend{hasil bagi} dan \definend{sisa}, yaitu \( q(x) \) dan \( r(x) \),
sedemikian sehingga
\begin{equation*}
  p(x)=d(x)\cdot q(x)+r(x)
\end{equation*}
dengan derajat \( r(x) \) lebih kecil secara ketat daripada derajat \( d(x) \).
%</th:EuclidForPolys>
\end{theorem}

Inti pernyataan bilangan bulat `\( 4 \) termuat \( 5 \) kali dalam \( 21 \)
dengan sisa \( 1 \)' adalah bahwa sisanya lebih kecil daripada \( 4 \):~walaupun
\( 4 \) termuat \( 5 \) kali, bilangan itu tidak termuat \( 6 \) kali.
Demikian pula, klausa terakhir dalam pernyataan pembagian polinomial sangat
penting.

\begin{example}
Jika \( p(x)=2x^3-3x^2+4x \) dan \( d(x)=x^2+1 \), maka \( q(x)=2x-3 \) dan
\( r(x)=2x+3 \).
Perhatikan bahwa derajat \( r(x) \) lebih kecil daripada derajat \( d(x) \).
\end{example}

\begin{corollary} \label{co:PolyDividedByLinearPolyIsConstant}
%<*co:PolyDividedByLinearPolyIsConstant>
Sisa pembagian \( p(x) \) oleh \( x-\lambda \) adalah polinomial konstan
\( r(x)=p(\lambda) \).
%</co:PolyDividedByLinearPolyIsConstant>
\end{corollary}

\begin{proof}
%<*pf:PolyDividedByLinearPolyIsConstant>
Sisa tersebut harus berupa polinomial konstan karena derajatnya lebih kecil
daripada derajat pembagi \( x-\lambda \).
Untuk menentukan konstantanya, ambil pembagi $d(x)$ dalam teorema sebagai
$x-\lambda$, lalu substitusikan \( \lambda\/ \) bagi $x$.
%</pf:PolyDividedByLinearPolyIsConstant>
\end{proof}

%<*PolynomialFactor>
Jika suatu pembagi \( d(x) \) membagi habis polinomial yang dibagi \( p(x) \),
yang berarti bahwa \( r(x) \) merupakan polinomial nol, maka \( d(x) \) disebut
faktor\index{faktor}\index{polinomial!faktor} dari \( p(x) \).
Setiap akar\index{akar}\index{polinomial!akar} faktor tersebut, yaitu
setiap \( \lambda\in\Re \) yang memenuhi \( d(\lambda)=0 \), juga merupakan akar
\( p(x) \), sebab \( p(\lambda)=d(\lambda)\cdot q(\lambda)=0 \).
%</PolynomialFactor>

\begin{corollary}  \label{co:RootOfPolyIsAssocLinearFactor}
%<*co:RootOfPolyIsAssocLinearFactor>
Jika \( \lambda \) merupakan akar polinomial \( p(x) \), maka \( x-\lambda \)
membagi habis \( p(x) \); artinya, $x-\lambda$ merupakan faktor $p(x)$.
%</co:RootOfPolyIsAssocLinearFactor>
\end{corollary}

\begin{proof}
%<*pf:RootOfPolyIsAssocLinearFactor>
Menurut korolari di atas, $p(x)=(x-\lambda)\cdot q(x)+p(\lambda)$.
Karena $\lambda$ merupakan akar, $p(\lambda)=0$, sehingga $x-\lambda$
merupakan faktor.
%</pf:RootOfPolyIsAssocLinearFactor>
\end{proof}

Akar berulang suatu polinomial adalah bilangan $\lambda$ sedemikian sehingga
polinomial tersebut habis dibagi oleh $(x-\lambda)^n$ untuk suatu pangkat yang
lebih besar daripada satu.
Pangkat terbesar seperti itu disebut \definend{multiplisitas}
\index{multiplisitas akar}\index{polinomial!multiplisitas} dari~$\lambda$.

Mencari akar dan faktor polinomial berderajat tinggi dapat menjadi sulit.
Namun, bagi polinomial berderajat dua kita memiliki rumus kuadrat:~akar-akar
\( ax^2+bx+c \) adalah
\begin{equation*}
   \lambda_1=\frac{-b+\sqrt{b^2-4ac}}{2a}
   \qquad
   \lambda_2=\frac{-b-\sqrt{b^2-4ac}}{2a}
\end{equation*}
(jika diskriminan \( b^2-4ac \) negatif, polinomial tersebut tidak memiliki
akar bilangan real).
Polinomial yang tidak dapat difaktorkan menjadi dua polinomial berderajat lebih
rendah dengan koefisien real disebut tak tereduksi di atas bilangan real.
\index{polinomial tak tereduksi}\index{polinomial!tak tereduksi}

\begin{theorem}  \label{th:CubicsAndHigherFactor}
%<*th:CubicsAndHigherFactor>
Setiap polinomial konstan atau linear tak tereduksi di atas bilangan real.
Suatu polinomial kuadrat tak tereduksi di atas bilangan real jika dan hanya jika
diskriminannya negatif.
Tidak ada polinomial kubik atau berderajat lebih tinggi yang tak tereduksi di
atas bilangan real.
%</th:CubicsAndHigherFactor>
\end{theorem}

\begin{corollary}  \label{co:RealPolysFactorIntoLinearsAndQuads}
%<*co:RealPolysFactorIntoLinearsAndQuads>
Setiap polinomial berkoefisien real dapat difaktorkan menjadi hasil kali
polinomial linear dan polinomial kuadrat tak tereduksi berkoefisien real.
Faktorisasi tersebut unik; sebarang dua faktorisasi memiliki faktor-faktor yang
sama dengan pangkat yang sama pula.
%</co:RealPolysFactorIntoLinearsAndQuads>
\end{corollary}

Perhatikan analoginya dengan faktorisasi prima bilangan bulat.
Dalam kedua kasus, klausa keunikan sangat berguna.

\begin{example}
Karena keunikan, tanpa mengalikan semua faktornya kita mengetahui bahwa
\( (x+3)^2(x^2+1)^3 \) tidak sama dengan \( (x+3)^4(x^2+x+1)^2 \).
\end{example}

\begin{example}
Menurut keunikan, jika \( c(x)=m(x)\cdot q(x) \), dengan
\( c(x)=(x-3)^2(x+2)^3 \) dan \( m(x)=(x-3)(x+2)^2 \), maka kita mengetahui
bahwa \( q(x)=(x-3)(x+2) \).
\end{example}

%<*ComplexNumbers>
Walaupun \( x^2+1 \) tidak memiliki akar real sehingga tidak dapat difaktorkan
di atas bilangan real, andaikan terdapat suatu akar\Dash yang secara tradisional
dilambangkan dengan \( i \)\Dash sedemikian sehingga \( i^2+1=0 \). Dengan
demikian, \( x^2+1 \) dapat difaktorkan menjadi hasil kali polinomial linear,
\( (x-i)(x+i) \).
Jika akar \( i \) ini kita tambahkan pada bilangan real dan sistem baru tersebut
kita tutup terhadap penjumlahan dan perkalian, kita memperoleh bilangan
kompleks,\index{bilangan kompleks}
\( \C=\set{a+bi\suchthat \text{$a,b\in\R$ dan $i^2=-1$}} \).
%</ComplexNumbers>

Bagi skalar $z\in\C$, dengan $z=a+bi$, kita menyebut $a$ sebagai
\definend{bagian riil} dari~$z$ dan~$b$ sebagai \definend{bagian imajiner}.
Bilangan kompleks sering kita gambarkan pada \definend{bidang kompleks}
\index{bidang kompleks}, dengan $a$ digambar pada \definend{sumbu riil}
\index{sumbu riil}, yaitu sumbu horizontal, dan $b$ digambar pada
\definend{sumbu imajiner}\index{sumbu imajiner}, yaitu sumbu vertikal.
Perhatikan bahwa jarak titik tersebut dari titik asal adalah panjangnya,
$|a+bi|=\sqrt{a^2+b^2}$.

%<*ComplexOperations>
Ingat kembali definisi penjumlahan bilangan kompleks dan perkalian skalar,
\begin{equation*}
   (a+bi)\,+\,(c+di)=(a+c)+(b+d)i                  
  \qquad
  r\cdot(a+bi)=(ra)+(rb)i
\end{equation*}
(sehingga pengurangannya adalah $(a+bi)-(c+di)=(a-c)+(b-d)i$).
Ingat pula definisi perkalian dua bilangan kompleks.
\begin{align*}
     (a+bi)(c+di) &=ac+adi+bci+bd(-1)  \\
                  &=(ac-bd)+(ad+bc)i
\end{align*}
%</ComplexOperations>

\begin{example}
Sebagai contoh,
% \( (1-2i)\,+\,(5+4i)=6+2i \) and
\( (2-3i)(4-0.5i)=6.5-13i \).
\end{example}

Di atas bilangan kompleks, setiap polinomial kuadrat dapat difaktorkan menjadi
polinomial-polinomial linear.
\begin{equation*}
   ax^2+bx+c=
   a\cdot \big(x-\frac{-b+\sqrt{b^2-4ac}}{2a}\big)
    \cdot \big(x-\frac{-b-\sqrt{b^2-4ac}}{2a}\big)
\end{equation*}

\begin{example}
Polinomial berderajat dua \( x^2+x+1 \) dapat difaktorkan di atas bilangan
kompleks menjadi hasil kali dua polinomial berderajat satu.
\begin{equation*}
   \big(x-\frac{-1+\sqrt{-3}}{2}\big)
   \big(x-\frac{-1-\sqrt{-3}}{2}\big)
   =
   \big(x-(-\frac{1}{2}+\frac{\sqrt{3}}{2}i)\big)
   \big(x-(-\frac{1}{2}-\frac{\sqrt{3}}{2}i)\big)
\end{equation*}
\end{example}

Berbeda dengan di atas bilangan real, dalam \( \C \) tidak ada polinomial
kuadrat yang tak tereduksi.
Semua polinomial dapat difaktorkan sepenuhnya menjadi polinomial linear.

\begin{theorem}[Teorema Dasar Aljabar] \label{th:FundThmAlg}\index{Teorema Dasar!Aljabar}
\hspace*{0em plus2em}
%<*th:FundThmAlg>
Polinomial berkoefisien kompleks dapat difaktorkan menjadi polinomial linear
berkoefisien kompleks.
Faktorisasi tersebut unik.
%</th:FundThmAlg>
\end{theorem}












\subsection{Representasi Kompleks}
Dengan definisi bilangan kompleks di atas, semua operasi yang telah kita
gunakan bagi ruang vektor real berlaku tanpa perubahan bagi ruang vektor dengan
skalar kompleks.

\begin{example}
Perkalian matriks tetap sama, walaupun perhitungannya dapat melibatkan lebih
banyak operasi aritmetika.
\begin{multline*}
  \begin{mat}
     1+1i  &2-0i  \\
      i    &-2+3i
  \end{mat}
  \begin{mat}
     1+0i  &1-0i  \\
     3i    &-i
  \end{mat}             \\
 \begin{aligned}
  &=\begin{mat}
     (1+1i)\cdot(1+0i)+(2-0i)\cdot(3i)   &(1+1i)\cdot(1-0i)+(2-0i)\cdot(-i) \\
     (i)\cdot(1+0i)+(-2+3i)\cdot(3i)     &(i)\cdot(1-0i)+(-2+3i)\cdot(-i)
  \end{mat}                                                  \\
  &=\begin{mat}
     1+7i  &1-1i  \\
    -9-5i  &3+3i
  \end{mat}
 \end{aligned}
\end{multline*}
\end{example}

Sebisa mungkin, kita akan menggunakan tanpa perubahan segala sesuatu dari
bab-bab sebelumnya.
Sebagai contoh, kita akan menyebut
\begin{equation*}
   \sequence{\colvec{1+0i \\ 0+0i \\ \vdots \\ 0+0i},
             \dots,
             \colvec{0+0i \\ 0+0i \\ \vdots \\ 1+0i}}
\end{equation*}
sebagai \definend{basis standar}\index{basis standar}\index{basis!standar}%
\index{basis standar!skalar bilangan kompleks}
bagi \( \C^n \) sebagai ruang vektor di atas $\C$ dan kembali melambangkannya
dengan \( \stdbasis_n \).
Sebagai contoh lain, $\polyspace_n$ akan menyatakan ruang vektor polinomial
berderajat~$n$ dengan koefisien kompleks.


































